\subsubsection{Classical vs. Quantum Correlation}

\begin{flushleft}
    \textcolor{Green3}{\faIcon{dice} \textbf{Classical Correlation}}
\end{flushleft}
A \definition{Classical Correlation} arises from \textbf{information} that \textbf{already exists before any measurement} is performed.

\highspace
In a classical system, correlated outcomes can be explained by assuming that the \textbf{system possesses predetermined properties}. These properties are often called \definition{Hidden Variables}, because they \hl{determine the outcome of future observations} even if the observers do not know their values.

\highspace
\begin{examplebox}[: Classical Correlation]
    For example, consider again two coins prepared so that they always show the same face. Before any observation takes place, the system is already in one of two possible configurations:
    \begin{itemize}
        \item Heads - Heads;
        \item Tails - Tails.
    \end{itemize}
    The measurement does not create the correlation. It merely reveals information that was already present.
\end{examplebox}

\highspace
Therefore, knowing the outcome observed by one observer immediately allows the other outcome to be inferred. No instantaneous influence is required, because the \textbf{correlation existed before the measurement}.

\highspace
\begin{flushleft}
    \textcolor{Green3}{\faIcon{atom} \textbf{Quantum Correlation}}
\end{flushleft}
A \definition{Quantum Correlation} arises from \textbf{measurements performed on an entangled quantum system}.

\highspace
As in the classical case, each observer sees random outcomes individually. However, the observed \textbf{correlations} possess a remarkable property: they \textbf{depend on the measurement basis chosen by the observers}.

\highspace
\begin{examplebox}[: Quantum Correlation]
    For example, if Alice and Bob share the Bell state
    \begin{equation*}
        \ket{\Phi^+}
        =
        \frac{\ket{00}+\ket{11}}{\sqrt2}
    \end{equation*}
    And both measure in the same basis, their outcomes are perfectly correlated. However, if they choose different measurement bases, the observed correlations change.
\end{examplebox}

\highspace
\begin{flushleft}
    \textcolor{DarkOrange3}{\faIcon{balance-scale} \textbf{The Main Difference}}
\end{flushleft}
In a \hl{classical system}, changing how an observer looks at the system does not alter the underlying correlation. The \hl{correlation is determined entirely by the hidden variables} that existed \textbf{before the measurement}.

\highspace
In contrast, \hl{quantum correlations} depend on the \hl{measurements that are actually performed}. The observed correlations cannot always be explained by assuming that the particles carried predetermined answers before the measurement.

\highspace
This distinction is what makes quantum correlation fundamentally different from classical correlation and lies at the heart of the EPR debate.

\highspace
\begin{flushleft}
    \textcolor{Green3}{\faIcon{book} \textbf{Bell Pair Example}}
\end{flushleft}
The distinction between classical and quantum correlation becomes clearer when we analyze an entangled Bell pair. Consider the Bell state:
\begin{equation*}
    \ket{\Phi^+}
    =
    \frac{\ket{00}+\ket{11}}{\sqrt{2}}
\end{equation*}
This state can also be expressed in the Hadamard basis as (see the full derivation in \autoref{example:bell-states-in-hadamard-basis}):
\begin{equation*}
    \ket{\Phi^+}
    =
    \frac{\ket{++}+\ket{--}}{\sqrt{2}}
\end{equation*}
Suppose Alice and Bob each possess one qubit of this entangled pair:
\begin{equation*}
    \ket{\Phi^+}_{AB}
    =
    \frac{\ket{00}_{AB}+\ket{11}_{AB}}{\sqrt{2}}
    =
    \underbrace{\frac{\ket{+}_{A} \otimes \ket{+}_{B}+\ket{-}_{A} \otimes \ket{-}_{B}}{\sqrt{2}}}_{\frac{\ket{0}_{A} \otimes \ket{0}_{B} + \ket{1}_{A} \otimes \ket{1}_{B}}{\sqrt{2}}}
\end{equation*}

\begin{examplebox}[: Same Basis Measurements]
    Suppose Alice and Bob both measure in the computational basis $\{\ket{0}, \, \ket{1}\}$.

    \highspace
    Alice measures first and obtains:
    \begin{equation*}
        \ket{0}
        \qquad\text{or}\qquad
        \ket{1}
    \end{equation*}
    With equal probability:
    \begin{equation*}
        P(0)=P(1)=\frac{1}{2}
    \end{equation*}
    If Alice obtains:
    \begin{equation*}
        \ket{0}
    \end{equation*}
    The Bell state collapses to:
    \begin{equation*}
        \ket{00}
    \end{equation*}
    Therefore Bob's qubit becomes:
    \begin{equation*}
        \ket{0}
    \end{equation*}
    Similarly, if Alice obtains:
    \begin{equation*}
        \ket{1}
    \end{equation*}
    The Bell state collapses to:
    \begin{equation*}
        \ket{11}
    \end{equation*}
    And Bob's qubit becomes:
    \begin{equation*}
        \ket{1}
    \end{equation*}
    Consequently, \hl{Bob always obtains the same outcome as Alice}.

    \highspace
    Although each observer sees \textbf{random outcomes individually}, they discover \textbf{perfect correlation when they later compare their measurement results}.
\end{examplebox}

\highspace
\begin{examplebox}[: Different Basis Measurements]
    Now suppose Alice measures in the Hadamard basis:
    \begin{equation*}
        \left\{
        \ket{+}, \, \ket{-}
        \right\}
    \end{equation*}
    Where
    \begin{equation*}
        \ket{+}
        =
        \frac{\ket{0}+\ket{1}}{\sqrt{2}},
        \qquad
        \ket{-}
        =
        \frac{\ket{0}-\ket{1}}{\sqrt{2}}
    \end{equation*}
    Bob, however, continues measuring in the computational basis.

    \highspace
    Alice obtains either:
    \begin{equation*}
        \ket{+}
        \qquad\text{or}\qquad
        \ket{-}
    \end{equation*}
    With probability:
    \begin{equation*}
        \frac{1}{2}
    \end{equation*}
    If Alice measures $\ket{+}$, Bob's qubit collapses to:
    \begin{equation*}
        \ket{+}
    \end{equation*}
    If Alice measures $\ket{-}$, Bob's qubit collapses to:
    \begin{equation*}
        \ket{-}
    \end{equation*}
    However, \hl{Bob is measuring in the computational basis}. For both $\ket{+}$ and $\ket{-}$:
    \begin{equation*}
        P(0)=P(1)=\frac{1}{2}
    \end{equation*}
    Therefore Bob obtains:
    \begin{equation*}
        0
        \qquad\text{or}\qquad
        1
    \end{equation*}
    With equal probability, regardless of Alice's outcome.

    \highspace
    When Alice and Bob later compare their results, they find that the \textbf{perfect correlation} observed in the previous experiment has \textbf{disappeared}.
\end{examplebox}

\highspace
\textcolor{Green3}{\faIcon[regular]{lightbulb}} In a classical hidden-variable model, the correlation would be determined entirely by pre-existing information carried by the particles. In contrast, the Bell-pair example shows that the \textbf{observed quantum correlation depends on the measurement basis chosen by the observers} (same basis, perfect correlation; different basis, no correlation). This basis dependence is one of the fundamental signatures of quantum correlation.